\documentclass[PoliticalScience]{vita}

\usepackage{color,hyperref}
\definecolor{darkblue}{rgb}{0.0,0.0,0.3}
\hypersetup{colorlinks,breaklinks,
            linkcolor=darkblue,urlcolor=darkblue,
            anchorcolor=darkblue,citecolor=darkblue}


\textheight 9in
\textwidth 6.5in
\topmargin -.5in
\oddsidemargin 0in
\evensidemargin 0in

\begin{document}

\name{Chad Hazlett  \\ \medskip Updated:  April 2026}

%\businessAddress{\small 
%Department of Statistics \\
%University of California, Los Angeles\\
%                 \small 8105G Math Sciences, Los Angeles, CA 90095-1554
% \\ \href{http://www.chadhazlett.com}{http://www.chadhazlett.com} \\
%                 \texttt{chazlett@ucla.edu}
%                 }%

\begin{vita}

\begin{Academic Positions}
\item Professor, Department of Statistics and Data Science, UCLA, July 2023-present
\item Professor, Department of Political Science, UCLA, July 2023-present
\item Associate Professor, Department of Statistics and Data Science, UCLA, July 2021-present
\item Associate Professor, Departments of Political Science, UCLA, July 2021-present
\item Assistant Professor, Department of Statistics and Data Science, UCLA, July 2014- June 2021
\item Assistant Professor, Departments of Political Science, UCLA, July 2014- June 2021
\end{Academic Positions}

\begin{Education}
  \item Ph.D., Political Science, Massachusetts Institute of Technology, 2009-2014.\\
  Dissertation:  \emph{Inference in Tough Places: Essays on Modeling, Matching, and Measurement with Applications to Civil Conflict.}\\
  Committee: Jens Hainmueller (chair), Fotini Christia, Adam Berinsky, Teppei Yamamoto\\
  Subfields: Political Methodology and International Relations

\medskip
\item Pre-doctoral Fellow (with Kosuke Imai), Dept. of Politics, Princeton, 2013-2014

\medskip
  \item M.P.P., Harvard University, Kennedy School of Government, 2004-2006. \\
  Concentration: \emph{Political and Economic Development}
\medskip
  \item B.S. with Distinction, \textit{magna cum laude}, Duke University, 1998--2002.
   \\ Major: Psychology; Concentration in Cognitive Neuroscience

\end{Education}

\medskip

\begin{Research Interests}
\item  Causal inference with observational data, non-randomized trials, sensitivity analysis, machine learning, measurement of sensitive attitudes and emotions, civil conflict, mass atrocities and genocide.
\end{Research Interests}

\medskip

\begin{Work in Progress}

\item Rohde, A., \& Hazlett, C. Causal progress with imperfect placebo treatments and outcomes. arXiv preprint arXiv:2310.15266. \textit{Revise and Resubmit, J. of Royal Statistical Society: A.}

\item Hazlett, C., \& Xu, Y. Trajectory balancing: A general weighting approach to causal inference with panel data. Available at SSRN 3214231.

%\item \textit{Under review}

%\item Wulf, D., Hazlett, C., Hill, B., Chiang, J., Goodman-Meza, D., Pasanuic, B., Arah, O., Erlandson, K., Montague, B. Safe learning outside of randomized trials: Application of the stability-controlled quasi-experiment to the effects of three COVID-19 therapies. \textit{Revise and Resubmit, Observational Studies}.

 
%\item Hartman, E., Hazlett, C., Sterbenz, C. kpop:  A kernel balancing approach for reducing specification assumptions in survey weighting.  \textit{R\&R at schoJRSS:A} 

%\item  Timothy S Chang, Yi Ding, Malika K Freund, Ruth Johnson, Tommer Schwarz,  Julie M Yabu, Chad Hazlett,  Jeffrey N Chiang, David Ami Wulf, UCLA Precision Health Data Discovery Repository Working Group, Daniel H Geschwind, Manish J Butte, Bogdan Pasaniuc.  Pre-existing conditions in Hispanics/Latinxs that are COVID-19 risk factors. \textit{Revise and resubmit},  iScience. Available on request. 

%\item Hazlett, C., Wulf, D., Hill, B., Chiang, J.  What can we credibly learn about the effects of remdesivir on COVID-19 mortality using electronic health records from outside of randomized trials? The stability-controlled quasi-experiment. Available at \href{http://www.chadhazlett.com}{http://www.chadhazlett.com}

%\item \textit{In preparation}
%\normalsize
%
%\item Hazlett, C., Wulf, D., Arah, O., Pasanuic, B., Erlandson, K., Montague, B.  Credible learning of hydroxychloroquine and dexamethasone effects on COVID-19 mortality outside of randomized trials. Available at \href{http://www.chadhazlett.com}{http://www.chadhazlett.com}
%
%
%\item Hazlett, C., Xu. Y. Trajectory Balancing: A general reweighting approach to causal inference with time-series cross-sectional data. %\textit{In preparation.}  
%Prior version available at \href{http://www.chadhazlett.com}{http://www.chadhazlett.com}
%
%%\item Inoue, K., Arah, O., Seaman, M., Cinelli, C., Hazlett C. Do prescription opioids increase all-cause mortality among the US general population? 
%
%\item Hazlett, C., Mello, F. What good is a free house? Evidence from 3000 housing lotteries in Brazil. Pre-analysis plan registration 20200219AB, available at \href{https://osf.io/yz8wa}{https://osf.io/yz8wa}. 
%
%%\item Hazlett, C., Sonnet, L.  Beyond Logit, Probit, or Linear Probability Models: avoiding misspecification bias while maintaining interpretability for binary outcome regressions.
%
%%\item Hazlett, C., Wulf, D.A., Hill, B., Chiang, J. What we can learn about Remdesivir from observational data: A stability-controlled quasi-experiment in a Los Angeles Health System.
%
%%\item Amini, A., Hazlett, C., Zhang L. Practical kernel regularized least squares with a million observations. 
%
%% \item Hazlett, C., Ramos, A., Smith, S. Using machine learning approaches to improve practical targeting of infant mortality prevention measures. 
% 
\end{Work in Progress}

\medskip

\begin{Publications}

\item \emph{Peer reviewed (authorship usually alphabetical)}

\item Hazlett, C., McMurry, N., Shinkre, T. (\textit{2026+}). Post-treatment problems: What can we say about the effect of a treatment among sub-groups who (would) respond in some way?  \textit{Forthcoming in Political Analysis}.  

\item Cho, S., Kim, D., \& Hazlett, C. (2026). Inference at the Data’s Edge: Gaussian Processes for Estimation and Inference in the Face of Extrapolation Uncertainty. \textit{Political Analysis}, 1-20.

\item Rapp, A. M., Ponting, C., Ramos, G., Escovar, E., Hazlett, C., Tan, P. Z., Torres, V. \& Chavira, D. A. (2026). A data-driven approach to identifying determinants of depression in rural Latine adolescents. \textit{Cultural Diversity \& Ethnic Minority Psychology}.

\item Hartman, E., Hazlett, C., Sterbenz, C. (2025). Kpop: A kernel balancing approach for reducing specification assumptions in survey weighting. \textit{Journal of the Royal Statistical Society Series A: Statistics in Society}. \href{https://doi.org/10.1093/jrsssa/qnae082}{https://doi.org/10.1093/jrsssa/qnae082}

\item Bertoli, A., \& Hazlett, C. (2025). Seeing like a district: Understanding what close-election designs for leader characteristics can and cannot tell us. \textit{Political Analysis}, 1-19.

\item Hazlett*, C., Wulf*, D.,  Hill, B., Chiang, J., Goodman-Meza, D., Pasanuic, B., Arah, O., Erlandson, K., Montague, B. (2025). Safe learning outside of randomized trials: Application of the stability-controlled quasi-experiment to the effects of three COVID-19 therapies. \textit{Observational Studies}. (Hazlett and Wulf co-first authors)

\item Cinelli, C., \& Hazlett, C. (2025). An omitted variable bias framework for sensitivity analysis of instrumental variables. \textit{Biometrika}, 112(2).

\item Cinelli, C., Ferwerda, J., Hazlett, C. (2024). sensemakr: Sensitivity Analysis Tools for OLS in R and Stata. \textit{Forthcoming, Observational Studies}. %Available at \href{http://www.chadhazlett.com}{http://www.chadhazlett.com}

\item Hazlett, C., Ramos, A., Smith, S. (2023). Better individual-level risk models can improve the targeting and life-saving potential of early-mortality interventions. \textit{Scientific Reports}, 13, 21706.

\item Faries, D., Gao, C., Zhang, X., Hazlett, C., Stamey, J., Yang, S., Ding, P., Shan, M., Sheffield, K., Dreyer, N. (2025). Real Effect or Bias? Best Practices for Evaluating the Robustness of Real-World Evidence through Quantitative Sensitivity Analysis for Unmeasured Confounding. \textit{Pharmaceutical Statistics 24, no. 2: e2457.}

\item Fabbe, K., Hazlett, C., Sinmazdemir, T. (2024). Threat Perceptions, Loyalties and Attitudes Towards Peace: The Effects of Civilian Victimization among Syrian Refugees in Turkey. \textit{Conflict Management and Peace Sciences}. %Available at \href{http://www.chadhazlett.com}{http://www.chadhazlett.com}



\item Cesar B. Martinez-Alvarez, Chad Hazlett, Paasha Mahdavi, and Michael L. Ross. (2023). Reply to van den Bergh and Savin: Fossil fuel taxes are politically hard to change. \textit{Proceedings of the National Academy of Sciences}, 2023

\item Hazlett, C., Parente F. (2023) From ``Is it unconfounded?''  to ``How much confounding would it take?'' A sensitivity-based approach to observational studies. (2023) \textit{Journal of Politics}. %\\ Available at \href{http://www.chadhazlett.com}{http://www.chadhazlett.com}

\item Cesar B. Martinez-Alvarez, Chad Hazlett, Paasha Mahdavi, and Michael L. Ross. Political Leadership Has Limited Impact on Fossil Fuel Taxes and Subsidies. (2022) \textit{Proceedings of the National Academy of Sciences}, 119 (47).

\item A. Akhazhanov, A. More, A. Amini, C. Hazlett, T. Treu, S. Birrer, and the Dark Energy Survey Collaboration. (2022). Finding quadruply imaged quasars with machine learning. I. Methods. \textit{Monthly Notices of the Royal Astronomical Society}, 513(2).

\item Graeme Blair, Mohammed Bukar, Rebecca Littman, Elizabeth R. Nugent, Rebecca Wolfe, Benjamin Crisman, Anthony Etim, Chad Hazlett, Jiyoung Kim. (2021). Trusted authorities can change minds and shift norms during conflict. \textit{Proceedings of the National Academcy of Science}, 118 (42). 

\item Chang, T.S., Ding, Y., Freund, M.K., Johnson, R., Schwarz, T., Yabu, J.M., Hazlett, C., Chiang, J.N., Wulf, D.A., Antonio, A.L. and Ariannejad, M. (2021). Pre-existing conditions in Hispanics/Latinxs that are COVID-19 risk factors. Iscience, 24(3), p.102188.

\item Blum, A., Hazlett, C., \& Posner, D. N. (2021). Measuring Ethnic Bias: Can Misattribution-Based Tools from Social Psychology Reveal Group Biases that Economics Games Cannot? \textit{Political Analysis}. 29(3), 385-404. %Available at \href{http://www.chadhazlett.com}{http://www.chadhazlett.com}

\item Conley, B., \& Hazlett, C. (2021). How very massive atrocities end: A dataset and typology.  \textit{Journal of Peace Research}, 58(3).  Data, case studies, and additional materials available at \href{http://www.chadhazlett.com}{http://www.chadhazlett.com}. 

\item Hazlett, C., Wainstein, L.  (2020). Understanding, choosing, and unifying multilevel and fixed effect approaches. Forthcoming, \textit{Political Analysis}: 1-20. %Available at \href{http://www.chadhazlett.com}{http://www.chadhazlett.com}

\item Hazlett, C., Maokola, W., Wulf, D. (2020). Inference without randomization or ignorability:  A stability-controlled quasi-experiment on the prevention of tuberculosis. \textit{Statistics in Medicine}, 39(28), 4149-4186. DOI: 10.1002/sim.8717. %Available at \href{http://www.chadhazlett.com}{http://www.chadhazlett.com}

\item Hazlett, C., \& Mildenberger, M. (2020). Wildfire exposure increases pro-climate political behaviors. \textit{American Political Science Review}. Online 15 July 2020. \href{doi:10.1017/S0003055420000441}{doi:10.1017/S0003055420000441}. %Available at \href{http://www.chadhazlett.com}{http://www.chadhazlett.com}

\item Hazlett, C. (2020). Kernel Balancing: A flexible non-parametric weighting procedure for estimating causal effects. \textit{Statistica Sinica}. %Available at \href{http://www.chadhazlett.com}{http://www.chadhazlett.com}

\item Hazlett, C., Campos E., Tan, P., Truong, H., Loo, S., DiStefano, C., Jeste, S., \& Senturk, D. (2020). Principle ERP reduction and analysis: Estimating and using principle ERP waveforms underlying ERPs across tasks, subjects and electrodes. \textit{NeuroImage}, 212, 116630. (Co-first author with  E. Campos). %Available at \href{http://www.chadhazlett.com}{http://www.chadhazlett.com}.

\item Cinelli, C., \& Hazlett, C. (2020). Making sense of sensitivity: Extending omitted variable bias. \textit{Journal of the Royal Statistical Society: Series B (Statistical Methodology)}, 82(1), 39-67. %\href{http://www.chadhazlett.com}{http://www.chadhazlett.com}.
	
\item Hazlett, C. (2020). Angry or Weary? How violence impacts attitudes toward peace among Darfurian refugees. \textit{Journal of Conflict Resolution}, 64(5), 844-870.  %Available at \href{http://www.chadhazlett.com}{http://www.chadhazlett.com}

\item Hazlett, C. (2019). Estimating causal effects of new treatments despite self-selection: The case of experimental medical treatments. \textit{Journal of Causal Inference}, 7(1).  %Available at \href{http://www.chadhazlett.com}{http://www.chadhazlett.com}. 

\item Fabbe, K., Hazlett, C., \& Sinmazdemir, T. (2019). A persuasive peace: Syrian refugees’ attitudes towards compromise and civil war termination. \textit{Journal of Peace Research}, 56(1), 103-117. %``A Persuasive Peace: Syrian Refugees’ Attitudes towards Compromise and Civil War Termination'' (with Kristin Fabbe and Tolga Sinmazdemir). \textit{Journal of Peace Research} (2019). 
%Available at \href{http://www.chadhazlett.com}{http://www.chadhazlett.com}.
 
\item Fong, C., Hazlett, C., \& Imai, K. (2018). Covariate balancing propensity score for a continuous treatment: Application to the efficacy of political advertisements. \textit{The Annals of Applied Statistics}, 12(1), 156-177.  %``Covariate Balancing Propensity Score for a Continuous Treatment: Application to the Efficacy of Political Advertisements'' (with Christian Fong and Kosuke Imai).  2018. \textit{Annals of Applied Statistics}, 12(1), 156-177. 
%Available at \href{http://www.chadhazlett.com}{http://www.chadhazlett.com}.

\item  Hazlett, C., \& Berinsky, A. (2018). Stress-testing the affect misattribution procedure: Heterogeneous control of affect misattribution procedure effects under incentives. \textit{British Journal of Social Psychology}, 57(1), 61-74.
%``Stress-testing the Affect Misattribution Procedure'' (with Adam Berinsky). September 2017. \textit{British Journal of Social Psychology}. 
%Available at \href{http://www.chadhazlett.com}{http://www.chadhazlett.com}.

\item  Ross, M. L., Hazlett, C., \& Mahdavi, P. (2017). Global progress and backsliding on gasoline taxes and subsidies. \textit{Nature Energy}, 2(1), 1-6. %`Global progress and backsliding on gasoline taxes and subsidies'' (with Michael Ross \& Paasha Mahdavi). \textit{Nature Energy}, January 2017.

\item  Ferwerda, J., Hainmueller, J., \& Hazlett, C. (2017). Kernel-Based Regularized Least Squares in R (KRLS) and Stata (krls). \textit{Journal of Statistical Software}, 79(3).
%Available at \href{http://www.chadhazlett.com}{http://www.chadhazlett.com}.

\item de Waal, A., Davenport, C., Hazlett, C., Kennedy, J. (2014). The Epidemiology of Lethal Violence in Darfur: Using Micro-Data to Explore Complex Patterns of Ongoing Armed Conflict. \textit{Social Science \& Medicine}, 120.

\item  Hainmueller, J., Hazlett, C. (2014). Kernel Regularized Least Squares: Reducing Misspecification Bias with a Flexible and Interpretable Machine Learning Approach. \textit{Political Analysis}, 22(2).

\item \textit{Prior to PhD}

\item Weissman, D. H., Gopalakrishnan, A., Hazlett, C. J., \& Woldorff, M. G. (2005). Dorsal anterior cingulate cortex resolves conflict from distracting stimuli by boosting attention toward relevant events. \textit{Cerebral Cortex}, 15(2), 229-237.

\item  Hazlett, C., Woldorff, M.G. (2004). Mechanisms of Moving the Mind's Eye: Planning and Execution of Spatial Shifts of Attention. \emph{Journal of Cognitive Neuroscience}, 16(5).

\item Woldorff, M. G., Hazlett, C. J., Fichtenholtz, H. M., Weissman, D. H., Dale, A. M., \& Song, A. W. (2004). Functional parcellation of attentional control regions of the brain. \textit{Journal of Cognitive Neuroscience}, 16(1), 149-165.

\item  Weissman, D. H., Woldorff, M. G., Hazlett, C. J., \& Mangun, G. R. (2002). Effects of practice on executive control investigated with fMRI. \textit{Cognitive Brain Research}, 15(1), 47-60.

 \item \emph{Software}

\item Please see \textit{software} section of my website for updates.

    \item \texttt{sensemakr}: Examining the sensitivity of regression results to unobserved confounding.\\
    \begin{itemize}
    \item    \texttt{R} software on github: \href{https://github.com/chadhazlett/sensemakr}{https://github.com/chadhazlett/sensemakr}
    \item \texttt{R} software on CRAN:  Install using  \texttt{install.packages(sensemakr)}.
    \item \texttt{Stata} software: Install using \texttt{ssc install sensemakr}
    \end{itemize}
    
    \item \texttt{kbal}:  An \texttt{R} package for Kernel Balancing.
    
    \begin{itemize}
    \item Development version available at \href{https://github.com/chadhazlett/kbal}{https://github.com/chadhazlett/kbal}
    \item Install with \texttt{devtools::install\_github("chadhazlett/kbal")}
    \end{itemize}
        
    \item \texttt{KRLS}:  Kernel regularized least squares %An \texttt{R} package for Kernel Regularized Least Squares (with Jens Hainmueller). A statistical software that implements Kernel Regularized Least Squares. Also available for Stata (with Jeremy Ferwerda and Jens Hainmueller). \\ \href{http://www.inside-r.org/packages/cran/KRLS/docs/krls}{http://www.inside-r.org/packages/cran/KRLS/docs/krls}
        \begin{itemize}
    \item \texttt{R} software on CRAN:  Install using  \texttt{install.packages(KRLS)}.
    \item \texttt{Stata} software: Install using \texttt{ssc install krls}
    \end{itemize}  
    
   \item  \texttt{pERPred}. Software for Principle ERP reduction and analysis: Estimating and using principle ERP waveforms underlying ERPs across tasks, subjects and electrodes (Hazlett and Campos, 2020). 
   \begin{itemize}
   \item Available on GitHub: \texttt{  devtools::install\_github("emjcampos/pERPred")}
   \end{itemize}
 
    \item Stability-controlled quasi-experiment (SCQE):
    \begin{itemize}
    \item \texttt{scqe}: available at \href{https://github.com/chadhazlett/scqe}{https://github.com/chadhazlett/scqe}.
    \item Install with \texttt{devtools::install\_github("chadhazlett/scqe")}
    \item Browser-based application for the stability controlled quasi-experiment  with aggregate data \\  \href{https://amiwulf.shinyapps.io/SCQE\_demo/}{https://amiwulf.shinyapps.io/SCQE\_demo/}. 
   \end{itemize}
\medskip

\item \emph{Book Chapters}

\item ``Beset on all sides: Children and the Landscape of Conflict in
North East Nigeria.''  (With Hilary Matfess and  Graeme Blair). In Cradled by Conflict: Child Involvement with Armed Groups in Contemporary Conflict. Ed: Siobhan O’Neil, Kato Van Broeckhoven.

\item ``Not On Our Watch: American Mobilization for Darfur'' (with R.J. Hamilton). In \emph{War in Darfur and the Search for Peace}, Alex De Waal (Ed.). Cambridge, Ma. Harvard University Press. 2007.

\medskip
\item \emph{News Media}

\item ``What Do Syrians Want Their Future to Be? A Survey of Refugees in Turkey'' (with Kristin Fabbe and Tolga Sinmazdemir). \textit{Foreign Affairs}.  May, 2017.

\item ``What Syrian refugees tell us about their hopes and fears for the future'' (with Kristin Fabbe and Tolga Sinmazdemir). \textit{Monkey Cage, Washington Post}. March, 2016. 

\item ``Primed to Protest'' (with Nicholas Miller). \emph{Foreign Policy.com}. Sept, 2012.

\item  ``India's Other Terrorism''. \emph{Far Eastern Economic Review}. Dec, 2008.

\end{Publications}

\medskip

%\begin{Teaching Experience}
%\item Teaching Assistant for Short Course on Causal Inference, Texas A\&M. January 2014. Instructor: Teppei Yamamoto. Constructed and taught materials for laboratory component. (Scheduled)
%\item Teaching Assistant for Quantitative Research Methods IV. MIT Political Science, Spring 2013.  Instructors: Teppei Yamamoto \& Danny Hidalgo. Overall Rating: 7 (out of 7).
%\item Teaching Assistant for Quantitative Research Methods II: Advanced Empirical Methods. MIT Political Science, Spring 2012.  Instructor: Jens Hainmueller. Overall Rating: 6.8 (out of 7).
%\item Instructor for ``Math Camp''. MIT Political Science, August 2010 \& 2011. Developed all teaching materials, lecture notes, handouts, and problem sets as first instructor for this course.
%\end{Teaching Experience}
%
%\medskip

\begin{Courses}
\item \textit{Advanced Causal Inference} (PhD students in Statistics, social science departments, business school, 25 students). 2018, 2021, 2022.
\item \textit{Causal Inference for Social Science}. (PhD and MS students in Political Science, Statistics, Health Sciences, Education, Sociology, other social science departments, enrollments ranged from 15-45 students). 2015, 2016, 2017, 2018, 2020, 2021, 2022.
\item \textit{Causal Inference for Social Science}. (Masters of Applied Statistics students, 17.). 2020, 2022 (forthcoming)
\item \textit{Statistical Modeling and Learning: Maximum likelihood and machine learning}. (PhD and MS students from Statistics, Engineering, Computer Science, other departments, 30-45 students.) 2015, 2016, 2017, 2018, 2019, 2020, 2021, 2022.
\item \textit{Machine Learning for Social Science: Current Literature.} (Small PhD reading group) 2016.
\item \textit{Undergraduate Statistics for Political Science.} (90-170 students). 2016, 2019, 2020, 2021, 2022.
%\item \textit{Statistics departmental speakers series} (Fall 2017, Winter, Spring 2018)
\item \textit{Mass Violence and Civil War} (PhD students in Political Science, Sociology). 2017.
\end{Courses}
\begin{Grants}

\item Bedari Kidness Institute.  ``From bias to better angels: Harnessing social desirability and self-
presentational demands to promote cross-group kindness''. \$18,000  (with Dan Posner and Jennifer Hamilton)
\item East Africa Social Science Translation Collaborative, Center for effective global action (2020). ``A piece of paper? The benefits of a secondary school certificate in Tanzania''.  \$70,000. (with Christina Fille)
\item Faculty Research Grant, ``How well does preventing evictions prevent homelessness?: A natural experiment''. UCLA (2019). \$6184. 
\item Facebook Statistics for Improving Insights and Decisions (2019), ``KPop: Kernel-based population reweighting of surveys''. \$50,000.  (with Erin Hartman)
\item California Center for Population Research Seed Grant (2018-2019),  ``Estimating Real World Causal Effects under Self-Selection", \$13,000.
\item Research Enabling Grant, UCLA Academic Senate (2018-2019), \$1800
\item Center for Social Statistics, Seed Grant (2018-2019). ``Improving Kernel Methods for Social Science Research.'' \$9,279. (with Arash Amini).
\item Research Enabling Grant, UCLA Academic Senate (2017-2018), \$1,800
\item Center for Social Statistics, Seed Grant (2017-2018). ``Randomized trials with varying probability of treatment''. \$20,000
\item California Center for Population Research Seed Grant (2016). ``The effects of Boko Haram violence on attitudes towards religion, policy, and peace.'' \$17,200  
\item Hellman Fellowship (2016-2017). ``Can social psychological tools detect ethnic biases that behavioral economics games miss? An initial test in Nairobi, Kenya.'' \$30,000
\item Research Enabling Grant, UCLA Academic Senate (2016). \$2000.
\item JPAL Governance Initiative (2015). ``Ruling (by) the Airwaves: The effect of radio programming on attitudes toward government in South Sudan.'' \$49,850. (Declined due to conflict escalation and suspension of operations by implementing partner)
\end{Grants}
\medskip

\begin{Fellowships}
\item US Holocaust Museum Simon-Skjodt Center Fellowship.  Statistical  analyses and  annual forecasts of mass atrocity risk. 2017-Present. 
\item Ford International Political Science Fellow, 2009-2013
\item \emph{India: Seva Mandir}. 2003-2004. Under Hart fellowship from  Duke University. Conducted field research with Seva Mandir, an NGO in Rajasthan, India. Examined access to health care in rural areas; designed and launched health insurance program.
\end{Fellowships}
\medskip

\begin{Service}
	\item Graduate Student Mentorship Award (2017)
	\item Methods Field Chair, Political Science (2016-2018)
	\item Co-founder, Southern California Methods Meeting (2016- )
    \item Faculty Search Committees:
    \begin{itemize} 
    \item Political Science: comparative politics (2014)
    \item Statistics: joint statistics \& social science search (2014), core of statistics (2022)
    \end{itemize}
    \item Manuscript Referee:
   \emph{American Political Science Review, J. of Causal Inference, J. of Experimental Political Science, J. of Genocide Research, J. of Conflict Resolution, J. of Public Policy, Political Science Research Methods, J. of Politics, Nature (Scientific Reports), J. of American Statistical Association, International Organization, Political Analysis, Political Behavior, Biometrics, Comparative Political Science, Statistica Sinica}
\end{Service}

\begin{Professional Memberships}
	\item UCLA Bedari Kindness Institute, Faculty Associate, 2020
	\item Center for California Population Research, Faculty Associate, 2014 onward
	\item Evidence in Governance and Politics (EGAP), Member, 2018 onward
	\item Center for Effective Global Action (CEGA), Faculty Affiliate,  2019 onward
\end{Professional Memberships}


%\begin{Invited Talks, Conferences}
%\item Asian PolMeth,  Jan 2022.
%\item Causal Data Science Meeting, November 2020
%\item Association for Public Policy Analysis \& Management (APPAM), November 2020
%\item Empirical Studies of Conflict (ESOC), September 2020
%\item Society for Political Methodology, Annual Summer Meeting. July 2020
%\item Workshop on Regression Discontinuity in West Africa (CEGA), May 2020
%\item UCLA Biostatistics Seminar, December 2019
%\item Brain Research Institute Electrophysiology of Brain Dynamics Affinity Group (UCLA), November 2019
%\item Southern California Methods, September 2019
%\item  East African Social Science Translation Collaborative Summit, Nairobi, July 2019
%\item Society for Political Methodology, Annual Summer Meeting. July 2019
%\item UC Santa Barbara Political Science, Methods Seminar. July 2019
%\item MIT Political Science, Methods Seminar, July 2019
%\item Southern California Methods, September 2018 
%\item Society for Political Methodology, Annual Summer Meeting. July 2018 
%\item Empirical Studies of Conflict, May 2018
%\item Working Group in African Politics, May 2018
%\item Causal Inference from Neuroscience to Computer Science,  UCLA. September 2017
%\item Southern California Methods (co-founder), September 2017
%\item American Political Science Association, Annual Meeting. San Francisco, August 2017  
%\item Empirical Studies of Conflict, May 2017.
%\item Southern California Methods (co-founder), August 2016.
%\item UC Riverside, Statistics Seminar, May 2016.
%\item USC Marshall, Data Sciences and Operations Department Seminar, March 2016. 
%\item Biostatistics Seminar (UCLA), December 2015. 
%\item International Political Economy Society, November 2015.
%\item Political Economy Workshop, Chicago Harris School of Public Policy, October 2015.
%\item California Center for Population Research Seminar, October 2015.
%\item American Mathematical Society, Western Regional Meeting, October 2015.
%\item ``Big Data and Death'' workshop, University of Wisconsin. December 2014.
%\item First Annual Leila and Melville Straus 1960 Family Symposium on Early Warning for Mass Atrocities. Dartmouth, October 2014
%\item American Political Science Association, Annual Meeting. Chicago, August 2014
%\item Stanford Political Methodology Workshop. April 2015.
%\item Society for Political Methodology, Annual Summer Meeting. University of Virginia. July 2013 
%\item New Faces in Political Methodology Conference. Penn State,  April 2013
%\item Midwest Political Science Association, Annual Meeting. Chicago, 2013, 2014, 2105, 2016
%\item B-WGAPE (Boston Working Group on African Political Economy).  November 2012.
%\item Applied Statistics Seminar (Institute for Quantitative Social Science).  Harvard University,  October 2012.
%\item MIT Strategic Use of Force Working Group.  Sept. 2012; May 2013
%\item Annual Meeting of the Northeast Methodology Program (NEMP). New York, May 2012.
%\item Midwest Political Science Association, Annual Meeting. Chicago, April 2012
%\item Lecture for AM111 (Harvard), ``Applied Mathematics: Scientific Computing''. Feb, 2012.
%%\item MIT Graduate Student Works in Progress. Feb. 2012.
%\end{Invited Talks, Conferences}

%\medskip

\begin{Consultations and Field Experience}

\item \emph{USAID: Reconciliation, Reintegration, and Stabilization in Mali}. November 2020-present. With Dan Posner and Salma Mousa. Developing and evaluating interventions aimed at facilitating reintegration and reducing inter-group conflict in Mali.

\item \emph{Early Warning Project, U.S. Holocaust Museum}. 2016-present. Providing annual risk forecasts and additional analysis of countries at risk of mass killings. 

\item \emph{Nigeria: United Nations University}. (2016-2017).  Preventing and responding to Boko Haram recruitment of children. 

\item \emph{Political Instability Task Force}. 2012-2018. Providing statistical analyses and forecasting models on variety of projects relating to mass killings and civil wars, using traditional and machine learning approaches.

\item \emph{U.S. Institute of Peace}. 2011. Contracted to assess existing genocide forecasting model and recommend improvements. Report available from U.S. Holocaust Museum at \href{http://bit.ly/14MZbfI}{http://bit.ly/14MZbfI}
%
%\item \emph{Kenya: Research Assistance to Fotini Christia (MIT)}. 2011. Conducted fieldwork for initial inquiry into effects of Somali piracy on markets in Nairobi.  Interviewed negotiators, attorneys, and policy-makers involved in Somali piracy.

\item \emph{Chad: U.S. Department of State}. 2010. Field statistician in eastern Chad for ``Darfurian Voices'', a project surveying over 2000 refugees, civilian leaders, and rebel leaders in refugee camps.

\item \emph{Mass Atrocities Response Operations, Carr Center for Human Rights, Harvard}. 2010.
Tasked to produce report on the status of African countries and regional organizations on their preparedness to respond to mass atrocities.

\item \emph{Burma, Thailand: Genocide Intervention Network}. 2008, 2010. Consultant to develop and evaluate a radio-based civilian early warning system in eastern Burma.

\item \emph{African Union}. 2009. Grant from the African Union Department of Peace and Stability to track and analyze levels of violence in Darfur in 2008-2009 using a variety of data sources from both the peacekeeping mission and media sources. Analyzed patterns of violence to help settle regarding nature of ongoing conflict. See ``The Epidemiology of Lethal Violence in Darfur: Using Event-Data to Explore Complex Patterns of Ongoing Armed Conflict''.
%
%\item \emph{Sudan, Ethiopia: Genocide Intervention Network}. 2007, 2008. Small-scale survey of IDPs in both government and rebel held territories of Darfur to determine the security risks they face. Worked with local partner and African Union peacekeeping forces to develop programs to improve camp security through improving patrols and by reducing need for firewood collection.
%
%\item \emph{Zimbabwe: Zimbabwe Opportunities Industrialization Center}. 2005. Conducted survey of small and medium-scale businesses able to survive economic crisis, shortages, and hyper-inflation to determine characteristics of resilient businesses.

\end{Consultations and Field Experience}

\medskip

%\begin{Manuscripts Under Review}
%\item NA
%\end{Manuscripts Under Review}

%\medskip

%\begin{Awards and Honors}
%\end{Awards and Honors}

%\begin{Grants}
%\end{Grants}

%\begin{Professional Experience}

%\end{Professional Experience}


%\medskip

%
%\begin{Field Experience}
%
%%\item \vspace{-.1in} Zimbabwe (2005). Conducted survey of small and medium-scale businesses able to survive economic crisis, shortages, and hyper-inflation to determine characteristics of resilient businesses.
%%\item \vspace{-.1in} Sudan (2007, 2008). Small-scale survey of IDPs to determine the security risks they faced; worked with local partner and African Union peacekeeping forces to develop programs to improve camp security through improving patrols and by reducing need for firewood collection.
%\item \vspace{-.1in} Chad (2009). Statistician  with U.S. State Department funded project to
%\item \vspace{-.1in} Burma (2008, 2010)
%\item \vspace{-.1in} Kenya (2011)
%\end{Field Experience}

%\begin{Press}
%\end{Press}

%
%\begin{Conference Participation}

%\end{Conference Participation}


%References
%\medskip
%\pagebreak[2]\par \textbf{\small{References}}\par
%\nopagebreak\bigskip \setlength{\tabcolsep}{.1em}
%\begin{tabular}{l@{\ \ \ \ \ \ \ \ \ \ }l}
%{
%\begin{tabular}{l}
%Jens Hainmueller \\
%Associate Professor of Political Science\\
%Department of Political Science \\
%Massachusetts Institute of Technology \\
%77 Massachusetts Avenue \\
%Cambridge, MA \ 02139  \\
%(617) 324-1935 \\
%\emph{jhainm@mit.edu} \\
%\end{tabular}
%} & {
%\begin{tabular}{l}
%Fotini Christia \\
%Assistant Professor of Political Science\\
%Department of Political Science\\
%Massachusetts Institute of Technology \\
%77 Massachusetts Avenue \\
%Cambridge, MA \ 02139 \\
%(617) 324-5595 \\
%\emph{cfotini@mit.edu} \\
%\end{tabular}
%}
%\\ \\
%{
%\begin{tabular}{l}
%Adam Berinsky \\
%Professor of Political Science \\
%Department of Political Science\\
%Massachusetts Institute of Technology \\
%77 Massachusetts Avenue \\
%Cambridge, MA \ 02139 \\
%(617) 253-8190 \\
%\emph{berinsky@mit.edu} \\
%\end{tabular}
%} & {
%\begin{tabular}{l}
%Teppei Yamamoto \\
%Assistant Professor of Political Science \\
%Department of Political Science\\
%Massachusetts Institute of Technology \\
%77 Massachusetts Avenue \\
%Cambridge, MA \ 02139 \\
%(617) 253-6959\\
%\emph{teppei@mit.edu} \\
%\end{tabular}
%}
%\\ \\
%{
%%\begin{tabular}{l}
%%Benjamin Valentino\\
%%Associate Professor of Political Science\\
%%Department of Government\\
%%Dartmouth College \\
%%208 Silsby Hall  \\
%%HB 6108 \\
%%Hanover, New Hampshire \ 03755 \\
%%(603) 646-2555 \\
%%\emph{benjamin.a.valentino@dartmouth.edu} \\
%%\end{tabular}
%}
%\end{tabular}

\end{vita}
\end{document}
